\begin{table}[htbp]
    \caption{HonestDiD Sensitivity for High-education share}
    \label{tab:honestdid-pct_voters_high_ed}
    \centering
    \small
    \renewcommand{\arraystretch}{1.08}
    \begin{tabular*}{0.94\textwidth}{@{\extracolsep{\fill}}lcccc}
        \doubletoprule
        Restriction & Violation bound & Lower bound & Upper bound & Rejects zero \\
        \midrule
        Smoothness & 0.0 & 0.085 & 0.094 & Yes \\
        Smoothness & 0.5 & -0.416 & 0.548 & No \\
        Smoothness & 1.0 & -0.916 & 1.048 & No \\
        Smoothness & 1.5 & -1.415 & 1.548 & No \\
        Smoothness & 2.0 & -1.912 & 2.048 & No \\
        \midrule
        Relative magnitudes & 0.0 & 0.085 & 0.094 & Yes \\
        Relative magnitudes & 0.5 & 0.084 & 0.095 & Yes \\
        Relative magnitudes & 1.0 & 0.082 & 0.096 & Yes \\
        Relative magnitudes & 1.5 & 0.080 & 0.098 & Yes \\
        Relative magnitudes & 2.0 & 0.078 & 0.100 & Yes \\
        \midrule
        Breakdown $\bar{M}$ & \multicolumn{4}{c}{>2.0} \\
        \doublebottomrule
    \end{tabular*}
    \vspace{1.0em}
    \begin{minipage}{0.94\textwidth}
        {\footnotesize \textbf{Note:} The table reports Rambachan-Roth sensitivity intervals for the event-time $0$ effect using the baseline dynamic TWFE specification. The smoothness rows impose bounds on second differences in latent violations of parallel trends. The relative-magnitudes rows bound post-treatment deviations relative to the largest pre-treatment deviation. The breakdown value is the smallest relative-magnitude violation at which the 95\% identified set contains zero; entries above the search grid are reported as lower bounds. \par}
    \end{minipage}
\end{table}
